% GENERATED FILE. Edit canonical lessons, then run npm run generate:books.
% canonical-source-hash: fnv1a64:9c7e0f11df623d95
% canonical-lessons: MR-C292-thingi, MR-C292-bhadakne, MR-C292-jvala, MR-C292-rakh, MR-C292-agnisaman-dal

\chapter{A Fire Breaks Out}
\label{ch:mr-a-fire-breaks-out}
\hlchaptermodality{292}

\begin{chapteropening}
\textbf{By the end of this chapter:} \emph{I can describe a fire: a spark, flames, ash and the fire brigade.}

Complete the last lesson of chapter 292: I can describe a fire: a spark, flames, ash and the fire brigade.
\end{chapteropening}

\section[\texorpdfstring{\mr{ठिणगी} (ṭhiṇgī)}{ठिणगी (ṭhiṇgī)}]{\mr{ठिणगी} (ṭhiṇgī) — a spark}

\label{lesson:MR-C292-thingi}



\begin{quote}
Before the new one: say the Marathi for missing, then the Marathi for to comb through.
\end{quote}

\subsection*{You'll want to know: \mr{ठिणगी}}
\textbf{\mr{ठिणगी}} — \emph{ṭhiṇgī} — \textquotedblleft{}a spark\textquotedblright{}.

\begin{quote}
\mr{एक} \mr{ठिणगी} \mr{पडल्यास} \mr{सगळं} \mr{गवत} \mr{पेटेल}.
\end{quote}

\begin{quote}
\emph{ek ṭhiṇgī paḍalyās sagḷaṁ gavat peṭel.}
\end{quote}

\begin{quote}
Should a single spark fall, all the grass will catch fire.
\end{quote}

Feminine. Plural: \mr{ठिणग्या}, \emph{ṭhiṇagyā}.

\mr{पडल्यास}, should it fall, says \mr{जर} … \mr{पडली} \mr{तर} in one word; plural \mr{ठिणग्या}.

\begin{cousinweb}
\textbf{\mr{ठिणगी}} (\emph{ṭhiṇgī}) is the everyday Marathi word for a spark; its origin is not settled.
\end{cousinweb}

\begin{grammarlens}[title={What you've built}]
One more word for telling what happened.
\end{grammarlens}

\subsection*{Your turn}
\begin{itemize}
\raggedright
  \item \emph{Say it:} \emph{ṭhiṇgī}
  \item \emph{Say it:} \emph{ṭhiṇgī}, once more
  \item \emph{Say it:} \emph{ek ṭhiṇgī paḍalyās sagḷaṁ gavat peṭel.} — Should a single spark fall, all the grass will catch fire.
  \item \emph{Recall:} say the Marathi for missing, then the Marathi for to comb through, then say \emph{ṭhiṇgī} again
\end{itemize}

\begin{tcolorbox}[breakable,colback=teal!4,colframe=teal!35!black,title={Before you move on}]
What does \mr{ठिणगी} mean? (\textquotedblleft{}A spark\textquotedblright{}.) Say it once more.
\end{tcolorbox}
\section[\texorpdfstring{\mr{भडकणे} (bhaḍakṇe)}{भडकणे (bhaḍakṇe)}]{\mr{भडकणे} (bhaḍakṇe) — to flare up; to blaze up}

\label{lesson:MR-C292-bhadakne}



\begin{quote}
Before the new one: say the Marathi for to comb through, then the Marathi for a spark.
\end{quote}

\subsection*{You'll want to know: \mr{भडकणे}}
\textbf{\mr{भडकणे}} — \emph{bhaḍakṇe} — \textquotedblleft{}to flare up; to blaze up\textquotedblright{}.

\begin{quote}
\mr{तेल} \mr{पडल्यावर} \mr{आग} \mr{भडकली}.
\end{quote}

\begin{quote}
\emph{tel paḍalyāvar āg bhaḍaklī.}
\end{quote}

\begin{quote}
When oil fell on it, the fire flared up.
\end{quote}

Past (he …): \mr{तो} \mr{भडकला}, \emph{to bhaḍaklā}.

Said of fire and of tempers: \mr{तो} \mr{भडकला}, he flared up in anger.

\begin{cousinweb}
\textbf{\mr{भडकणे}} (\emph{bhaḍakṇe}) is inherited: CDIAL lists Marathi \emph{bhaḍakṇe}, 'to blaze, take fright', under \emph{*bhaṭ}, 'a sudden movement or noise' (CDIAL 9365).
\end{cousinweb}

\begin{grammarlens}[title={What you've built}]
One more word for telling what happened.
\end{grammarlens}

\subsection*{Your turn}
\begin{itemize}
\raggedright
  \item \emph{Say it:} \emph{bhaḍakṇe}
  \item \emph{Say it:} \emph{bhaḍakṇe}, once more
  \item \emph{Say it:} \emph{tel paḍalyāvar āg bhaḍaklī.} — When oil fell on it, the fire flared up.
  \item \emph{Recall:} say the Marathi for to comb through, then the Marathi for a spark, then say \emph{bhaḍakṇe} again
\end{itemize}

\begin{tcolorbox}[breakable,colback=teal!4,colframe=teal!35!black,title={Before you move on}]
What does \mr{भडकणे} mean? (\textquotedblleft{}To flare up; to blaze up\textquotedblright{}.) Say it once more.
\end{tcolorbox}
\section[\texorpdfstring{\mr{ज्वाळा} (jvāḷā)}{ज्वाळा (jvāḷā)}]{\mr{ज्वाळा} (jvāḷā) — a tongue of flame; a blaze}

\label{lesson:MR-C292-jvala}



\begin{quote}
Before the new one: say the Marathi for a spark, then the Marathi for to flare up.
\end{quote}

\subsection*{You'll want to know: \mr{ज्वाळा}}
\textbf{\mr{ज्वाळा}} — \emph{jvāḷā} — \textquotedblleft{}a tongue of flame; a blaze\textquotedblright{}.

\begin{quote}
\mr{इमारतीच्या} \mr{खिडक्यांतून} \mr{ज्वाळा} \mr{बाहेर} \mr{येत} \mr{होत्या}.
\end{quote}

\begin{quote}
\emph{imārtīcyā khiḍakyāntūn jvāḷā bāher yet hotyā.}
\end{quote}

\begin{quote}
Flames were coming out of the windows of the building.
\end{quote}

Feminine.

\mr{ज्वाळा} keeps its form in the plural.

\begin{cousinweb}
\textbf{\mr{ज्वाळा}} (\emph{jvāḷā}) is a tatsama, taken from Sanskrit \emph{jvālā}, 'a flame' (Monier-Williams), with Marathi \mr{ळ} for the Sanskrit l.
\end{cousinweb}

\begin{grammarlens}[title={What you've built}]
One more word for telling what happened.
\end{grammarlens}

\subsection*{Your turn}
\begin{itemize}
\raggedright
  \item \emph{Say it:} \emph{jvāḷā}
  \item \emph{Say it:} \emph{jvāḷā}, once more
  \item \emph{Say it:} \emph{imārtīcyā khiḍakyāntūn jvāḷā bāher yet hotyā.} — Flames were coming out of the windows of the building.
  \item \emph{Recall:} say the Marathi for a spark, then the Marathi for to flare up, then say \emph{jvāḷā} again
\end{itemize}

\begin{tcolorbox}[breakable,colback=teal!4,colframe=teal!35!black,title={Before you move on}]
What does \mr{ज्वाळा} mean? (\textquotedblleft{}A tongue of flame; a blaze\textquotedblright{}.) Say it once more.
\end{tcolorbox}
\section[\texorpdfstring{\mr{राख} (rākh)}{राख (rākh)}]{\mr{राख} (rākh) — ash; ashes}

\label{lesson:MR-C292-rakh}



\begin{quote}
Before the new one: say the Marathi for to flare up, then the Marathi for a tongue of flame.
\end{quote}

\subsection*{You'll want to know: \mr{राख}}
\textbf{\mr{राख}} — \emph{rākh} — \textquotedblleft{}ash; ashes\textquotedblright{}.

\begin{quote}
\mr{आग} \mr{विझल्यावर} \mr{तिथे} \mr{फक्त} \mr{राख} \mr{उरली}.
\end{quote}

\begin{quote}
\emph{āg vijhalyāvar tithe phakta rākh urlī.}
\end{quote}

\begin{quote}
When the fire went out, only ash was left there.
\end{quote}

Feminine.

\mr{राख} \mr{होणे}, to burn to ashes.

\begin{cousinweb}
\textbf{\mr{राख}} (\emph{rākh}) is a tadbhava, inherited from Sanskrit \emph{rakṣā}, 'ashes' (CDIAL 10552).
\end{cousinweb}

\begin{grammarlens}[title={What you've built}]
One more word for telling what happened.
\end{grammarlens}

\subsection*{Your turn}
\begin{itemize}
\raggedright
  \item \emph{Say it:} \emph{rākh}
  \item \emph{Say it:} \emph{rākh}, once more
  \item \emph{Say it:} \emph{āg vijhalyāvar tithe phakta rākh urlī.} — When the fire went out, only ash was left there.
  \item \emph{Recall:} say the Marathi for to flare up, then the Marathi for a tongue of flame, then say \emph{rākh} again
\end{itemize}

\begin{tcolorbox}[breakable,colback=teal!4,colframe=teal!35!black,title={Before you move on}]
What does \mr{राख} mean? (\textquotedblleft{}Ash; ashes\textquotedblright{}.) Say it once more.
\end{tcolorbox}
\section[\texorpdfstring{\mr{अग्निशमन} \mr{दल} (agniśaman dal)}{अग्निशमन दल (agniśaman dal)}]{\mr{अग्निशमन} \mr{दल} (agniśaman dal) — a fire brigade}

\label{lesson:MR-C292-agnisaman-dal}



\begin{quote}
Before the new one: say the Marathi for a tongue of flame, then the Marathi for ash.
\end{quote}

\subsection*{You'll want to know: \mr{अग्निशमन} \mr{दल}}
\textbf{\mr{अग्निशमन} \mr{दल}} — \emph{agniśaman dal} — \textquotedblleft{}a fire brigade\textquotedblright{}.

\begin{quote}
\mr{अग्निशमन} \mr{दलाच्या} \mr{गाड्या} \mr{दहा} \mr{मिनिटांत} \mr{आल्या}.
\end{quote}

\begin{quote}
\emph{agniśaman dalācyā gāḍyā dahā miniṭānt ālyā.}
\end{quote}

\begin{quote}
The fire brigade arrived within ten minutes.
\end{quote}

Neuter.

\mr{दल} is neuter, and so is the whole: \mr{अग्निशमन} \mr{दलाला} \mr{फोन} \mr{करणे}, to call the fire brigade.

\begin{cousinweb}
\textbf{\mr{अग्निशमन} \mr{दल}} (\emph{agniśaman dal}) is a modern compound built from Sanskrit \emph{agni}, 'fire', \emph{śamana}, 'quelling', and \emph{dala}, 'a troop' (Monier-Williams).
\end{cousinweb}

\begin{grammarlens}[title={What you've built}]
One more word for telling what happened.
\end{grammarlens}

\subsection*{Your turn}
\begin{itemize}
\raggedright
  \item \emph{Say it:} \emph{agniśaman dal}
  \item \emph{Say it:} \emph{agniśaman dal}, once more
  \item \emph{Say it:} \emph{agniśaman dalācyā gāḍyā dahā miniṭānt ālyā.} — The fire brigade arrived within ten minutes.
  \item \emph{Recall:} say the Marathi for a tongue of flame, then the Marathi for ash, then say \emph{agniśaman dal} again
\end{itemize}

\begin{tcolorbox}[breakable,colback=teal!4,colframe=teal!35!black,title={Before you move on}]
What does \mr{अग्निशमन} \mr{दल} mean? (\textquotedblleft{}A fire brigade\textquotedblright{}.) Say it once more.
\end{tcolorbox}
